% Partial boundary observation: relative homology kernel, minimal audit completion, and matroidal structure.

\paragraph{部分边界观测的核与审计补全}
在定理 \ref{thm:spg-boundary-godelization-holographic-dictionary} 的“全边界可见”情形下，顶维边界单射给出严格可逆的离散全息字典。
当仅能观测部分取向面（屏幕）时，信息缺口由一个相对同调群精确刻画，并且其秩同时控制最小审计补全成本与可达的 Gödel 型补账寄存器维数。

\begin{theorem}[可计算双向编码下的 Kolmogorov 复杂度守恒：Stokes--Gödel 全息完备性]\label{thm:spg-stokes-godel-algorithmic-holographic-completeness}
固定 \(m,n\ge 1\)，并令 \(\mathcal Q_m\) 为 \(I^n=[0,1]^n\) 的 \(m\)-级 dyadic 立方复形。
取系数域 \(\FF_2\)，并记顶维边界算子为
\[
\partial_n:\ C_n(\mathcal Q_m;\FF_2)\longrightarrow C_{n-1}(\mathcal Q_m;\FF_2).
\]
由定理 \ref{thm:spg-dyadic-cubical-boundary-injective} 及 \(H_n(\mathcal Q_m;\FF_2)=0\) 可知 \(\partial_n\) 为单射。
令 \(G:C_{n-1}(\mathcal Q_m;\FF_2)\to\NN\) 为任意可计算且在其像上可计算可逆的平方自由 Gödel 整数化（例如对每条取向 \((n{-}1)\)-面分配互异素数并取其乘积），并定义
\[
\Phi:=G\circ \partial_n:\ C_n(\mathcal Q_m;\FF_2)\longrightarrow \NN.
\]
取任意固定的前缀型 Kolmogorov 复杂度 \(K(\,\cdot\,|\,\cdot\,)\)（相对于某个固定的通用前缀机）。
则存在仅依赖于 \(n\) 与所选编码/译码机的常数 \(C>0\)（与 \(m\) 与输入无关），使得对任意 \(u\in C_n(\mathcal Q_m;\FF_2)\) 都有
\[
\bigl|K(u\mid n,m)-K(\Phi(u)\mid n,m)\bigr|\le C,
\qquad
K(u\mid \Phi(u),n,m)\le C.
\]
并且存在可计算泛函 \(V:\NN\to\QQ_{\ge 0}\) 使得
\[
V(\Phi(u))=\operatorname{Vol}(u),
\]
其中 \(\operatorname{Vol}(u)\) 表示由 \(u\) 的占据数给出的 dyadic 体积（每个 \(n\)-胞腔体积取 \(2^{-mn}\)）。
\leanverified{paper\_spg\_stokes\_godel\_algorithmic\_holographic\_completeness}
\end{theorem}
\begin{proof}
由于 \(\partial_n\) 可计算且 \(G\) 可计算，\(\Phi\) 可计算，故 \(K(\Phi(u)\mid n,m)\le K(u\mid n,m)+O(1)\)。
另一方面，\(\partial_n\) 在 \(\FF_2\) 上单射且 \(\mathcal Q_m\) 有限，故 \(\partial_n^{-1}\) 在 \(\mathrm{im}(\partial_n)\) 上可由穷举/线性代数实现为可计算函数；
又由假设 \(G\) 在其像上可计算可逆，故 \(\Phi^{-1}\) 在 \(\mathrm{im}(\Phi)\) 上可计算。
于是 \(u\) 可由 \(\Phi(u)\) 与参数 \((n,m)\) 可计算恢复，给出
\(
K(u\mid \Phi(u),n,m)=O(1)
\)
并推出 \(K(u\mid n,m)\le K(\Phi(u)\mid n,m)+O(1)\)。
合并两侧不等式即得复杂度差为 \(O(1)\)。
最后，由 \(\Phi(u)\) 可计算得到 \(u\)，而 \(\operatorname{Vol}(u)\) 由 \(u\) 的占据数可计算给出，故存在可计算泛函 \(V\) 满足所示等式。证毕。
\end{proof}

\begin{theorem}[部分边界观测的屏幕核与相对顶维同调]\label{thm:spg-partial-boundary-screen-kernel-relative-homology}
固定 \(m,n\ge 1\) 并沿用定理 \ref{thm:spg-dyadic-cubical-boundary-injective} 的 dyadic 立方复形 \(\mathcal Q_m\) 与链群 \(C_k(\mathcal Q_m;\ZZ)\) 记号。
令 \(\overrightarrow{\mathcal F_{m}^{n-1}}\) 为全部取向 \((n-1)\)-胞腔所成的标准基。
对任意 \(S\subset \overrightarrow{\mathcal F_{m}^{n-1}}\)，令
\[
\Pi_S:C_{n-1}(\mathcal Q_m;\ZZ)\longrightarrow \ZZ^{S}
\]
为对 \(S\) 坐标的投影，并定义部分边界观测算子
\[
f_S:=\Pi_S\circ\partial_n:\ C_n(\mathcal Q_m;\ZZ)\longrightarrow \ZZ^{S}.
\]
令 \(A_{\neg S}\subset \mathcal Q_m^{(n-1)}\) 为从 \((n-1)\)-骨架中删去 \(S\) 中那些 \((n-1)\)-胞腔后得到的子复形（保留全部 \(\le n-2\) 维胞腔，因而仍为子复形，并且 \(C_n(A_{\neg S};\ZZ)=0\)）。
则有自然同构
\[
\ker f_S\ \cong\ H_n(\mathcal Q_m,\ A_{\neg S};\ZZ).
\]
\leanverified{paper\_spg\_partial\_boundary\_screen\_kernel\_relative\_homology}
\end{theorem}
\begin{proof}
由于 \(C_n(A_{\neg S};\ZZ)=0\)，相对链群满足 \(C_n(\mathcal Q_m,A_{\neg S};\ZZ)=C_n(\mathcal Q_m;\ZZ)\)。
相对边界算子为
\[
\partial_n^{\mathrm{rel}}:C_n(\mathcal Q_m;\ZZ)\longrightarrow C_{n-1}(\mathcal Q_m;\ZZ)/C_{n-1}(A_{\neg S};\ZZ).
\]
因此 \(x\in C_n(\mathcal Q_m;\ZZ)\) 为相对 \(n\)-循环当且仅当 \(\partial_n x\in C_{n-1}(A_{\neg S};\ZZ)\)，等价于 \(\Pi_S(\partial_n x)=0\)，亦即 \(x\in\ker f_S\)。
从而 \(\ker f_S=Z_n(\mathcal Q_m,A_{\neg S};\ZZ)\)。
\par
另一方面，\(\mathcal Q_m\) 无 \((n+1)\)-胞腔，故 \(C_{n+1}(\mathcal Q_m,A_{\neg S};\ZZ)=0\) 并且 \(B_n(\mathcal Q_m,A_{\neg S};\ZZ)=\mathrm{im}\,\partial_{n+1}^{\mathrm{rel}}=0\)。
因此
\[
H_n(\mathcal Q_m,A_{\neg S};\ZZ)=Z_n(\mathcal Q_m,A_{\neg S};\ZZ)/B_n(\mathcal Q_m,A_{\neg S};\ZZ)\cong \ker f_S.
\]
证毕。
\end{proof}

\begin{theorem}[屏幕核的连通分量闭式]\label{thm:spg-screen-kernel-connected-components}
\leanverified{paper\_spg\_screen\_kernel\_connected\_components}
在定理 \ref{thm:spg-partial-boundary-screen-kernel-relative-homology} 的设定下，
令 \(\Gamma\) 为 \(\mathcal Q_m\) 的顶维胞腔邻接图并添加外部顶点 \(\infty\)：其顶点集为
\[
V(\Gamma):=\mathcal Q_m^{(n)}\ \cup\ \{\infty\},
\]
其中 \(\mathcal Q_m^{(n)}\) 表示全部 \(n\)-胞腔（按单元而非取向计）。
由屏幕 \(S\subset \overrightarrow{\mathcal F_{m}^{n-1}}\) 诱导边集 \(E_S\) 如下：
\begin{enumerate}
  \item 若 \(\vec f\in S\) 的底层 \((n-1)\)-面为内部共享面，且其两侧顶维胞腔为 \(u,v\in\mathcal Q_m^{(n)}\)，则令 \(\{u,v\}\in E_S\)；
  \item 若 \(\vec f\in S\) 的底层 \((n-1)\)-面为外边界面，且其唯一相邻顶维胞腔为 \(v\in\mathcal Q_m^{(n)}\)，则令 \(\{v,\infty\}\in E_S\)。
\end{enumerate}
记 \(\Gamma_S:=(V(\Gamma),E_S)\) 的连通分量数为 \(c(\Gamma_S)\)。
则有自然同构
\[
\ker f_S\ \cong\ \ZZ^{\,c(\Gamma_S)-1}.
\]
\leanverified{kernel\_dim\_eq\_components\_minus\_one}
\leanverified{free\_components\_count}
\leanverified{fiber\_card\_binary}
\leanverified{screen\_monotone\_kernel}
\leanverified{paper\_spg\_screen\_kernel\_connected\_components\_seeds}
\leanverified{paper\_spg\_screen\_kernel\_connected\_components\_package}
\end{theorem}
\begin{proof}
将 \(x\in C_n(\mathcal Q_m;\ZZ)\) 写为 \(x=\sum_{v\in\mathcal Q_m^{(n)}}x(v)[v]\)，其中 \(x(v)\in\ZZ\) 为对应顶维胞腔系数（取一组固定的顶维取向基）。
对任意 \(\vec f\in S\)，\((\partial_n x)(\vec f)\) 仅依赖于与 \(\vec f\) 相邻的一个或两个顶维系数：
若 \(\vec f\) 为内部共享面且两侧为 \(u,v\)，则
\(
(\partial_n x)(\vec f)=\pm(x(u)-x(v))
\)；
若 \(\vec f\) 为外边界面且相邻胞腔为 \(v\)，则
\(
(\partial_n x)(\vec f)=\pm x(v)
\)。
因此 \(f_S(x)=0\) 当且仅当满足如下约束：
\[
u\text{ 与 }v\text{ 由 }E_S\text{ 相连}\ \Longrightarrow\ x(u)=x(v),\qquad
v\text{ 与 }\infty\text{ 由 }E_S\text{ 相连}\ \Longrightarrow\ x(v)=0.
\]
从而 \(x(\cdot)\) 在 \(\Gamma_S\) 的每个连通分量上取常值，且包含 \(\infty\) 的分量常值被迫为 \(0\)。
令 \(\Phi:\ker f_S\to \ZZ^{c(\Gamma_S)-1}\) 为“读取其余各分量常值”的映射，则由上式刻画知 \(\Phi\) 为双射群同态，故为同构。证毕。
\end{proof}

\begin{corollary}[部分屏幕的唯一回放判据与最小审计补全成本闭式]\label{cor:spg-screen-injectivity-and-audit-cost-components}
在定理 \ref{thm:spg-screen-kernel-connected-components} 的设定下，以下条件等价：
\begin{enumerate}
  \item \(f_S\) 为单射；
  \item \(c(\Gamma_S)=1\)（亦即 \(\Gamma_S\) 连通）。
\end{enumerate}
并且最小审计补全成本满足
\[
\min_{(R,r_S)}\cdim(R)
=\rank_{\ZZ}(\ker f_S)
=c(\Gamma_S)-1.
\]
\leanverified{paper\_spg\_screen\_injectivity\_and\_audit\_cost\_components}
\end{corollary}
\begin{proof}
由定理 \ref{thm:spg-screen-kernel-connected-components}，
\(
\ker f_S\cong \ZZ^{c(\Gamma_S)-1}
\)。
故 \(f_S\) 单射当且仅当 \(\ker f_S=0\)，等价于 \(c(\Gamma_S)=1\)，得到前半部。
最小审计补全成本等于 \(\rank_{\ZZ}(\ker f_S)\) 由推论 \ref{cor:spg-partial-boundary-min-audit-cost-kernel-rank} 给出；
代入闭式即得结论。证毕。
\end{proof}

\begin{proposition}[轴向屏幕的面积律：审计成本的指数级下界可达]\label{prop:spg-axial-screen-area-law-audit-cost}
固定坐标方向 \(x_1\)。
令 \(S^{(1)}\subset \overrightarrow{\mathcal F_{m}^{n-1}}\) 为所有\emph{正交于 \(x_1\)} 的内部 \((n-1)\)-面（即两顶维胞腔仅沿 \(x_1\) 方向相邻时所共享的面；不含任何外边界面）。
则
\[
c(\Gamma_{S^{(1)}})=2^{m(n-1)}+1,
\qquad
\min_{(R,r_{S^{(1)}})}\cdim(R)=2^{m(n-1)}.
\]
\end{proposition}
\leanverified{paper\_spg\_axial\_screen\_area\_law\_audit\_cost}
\begin{proof}
在 \(\Gamma_{S^{(1)}}\) 中，两顶维胞腔相连当且仅当其 dyadic 坐标仅在 \(x_1\) 方向相邻。
因此对每个固定的 \((i_2,\dots,i_n)\in\{0,\dots,2^m-1\}^{n-1}\)，全体
\(
(i_1,i_2,\dots,i_n)
\)
（\(i_1\) 变动）构成一条长度 \(2^m\) 的路径图连通分量；不同 \((i_2,\dots,i_n)\) 之间无边相连。
故由内部面诱导的连通分量共有 \(2^{m(n-1)}\) 个。
又因为 \(S^{(1)}\) 不含外边界面，顶点 \(\infty\) 为孤立分量，从而
\(
c(\Gamma_{S^{(1)}})=2^{m(n-1)}+1
\)。
最后由推论 \ref{cor:spg-screen-injectivity-and-audit-cost-components} 得最小审计补全成本为 \(c(\Gamma_{S^{(1)}})-1=2^{m(n-1)}\)。证毕。
\end{proof}

\begin{corollary}[\texorpdfstring{$\FF_2$}{F2} 计数下的纤维基数：相对 Betti 数的精确公式]\label{cor:spg-partial-boundary-fiber-cardinality-f2}
在定理 \ref{thm:spg-partial-boundary-screen-kernel-relative-homology} 的设定下，将系数域改为 \(\FF_2\)，并把 \(f_S\) 视为线性映射
\[
f_S:\ C_n(\mathcal Q_m;\FF_2)\longrightarrow \FF_2^{S}.
\]
对任意 \(y\in \mathrm{im}(f_S)\)，其纤维基数满足精确恒等式
\[
\card{f_S^{-1}(y)}=2^{\beta_n(\mathcal Q_m,A_{\neg S};\FF_2)},
\qquad
\beta_n(\mathcal Q_m,A_{\neg S};\FF_2):=\dim_{\FF_2}H_n(\mathcal Q_m,A_{\neg S};\FF_2).
\]
特别地，上式右端与 \(y\) 无关，仅依赖于屏幕 \(S\)。
\leanverified{paper\_spg\_partial\_boundary\_fiber\_cardinality\_f2}
\end{corollary}
\begin{proof}
由于 \(f_S\) 为 \(\FF_2\)-线性映射，任意非空纤维 \(f_S^{-1}(y)\) 为仿射子空间，故
\(
\card{f_S^{-1}(y)}=2^{\dim_{\FF_2}\ker f_S}
\)。
由定理 \ref{thm:spg-partial-boundary-screen-kernel-relative-homology} 的论证在 \(\FF_2\) 系数下同样成立，并且由于 \(\mathcal Q_m\) 无 \((n+1)\)-胞腔，仍有
\(
H_n(\mathcal Q_m,A_{\neg S};\FF_2)\cong \ker f_S
\)。
代入即得结论。证毕。
\end{proof}

\begin{corollary}[最小审计补全成本等于屏幕核秩]\label{cor:spg-partial-boundary-min-audit-cost-kernel-rank}
沿用定理 \ref{thm:spg-partial-boundary-screen-kernel-relative-homology} 的记号，并令 \(r:=\rank_{\ZZ}(\ker f_S)\)。
则任意使 \(f_S\) 可单射提升的审计补全 \((R,r_S)\)（定义 \ref{def:cdim-audit-completion-loss}）都满足
\[
\cdim(R)\ \ge\ r,
\]
并且该下界可达。
因此，\(f_S\) 的最小圆维审计补全成本等于
\[
\min_{(R,r_S)}\cdim(R)=\rank_{\ZZ}(\ker f_S)=\rank_{\ZZ}H_n(\mathcal Q_m,A_{\neg S};\ZZ).
\]
\leanverified{paper\_spg\_partial\_boundary\_min\_audit\_cost\_kernel\_rank}
\end{corollary}
\begin{proof}
由于 \(C_n(\mathcal Q_m;\ZZ)\) 与 \(\ZZ^{S}\) 均为自由阿贝尔群，\(\ker f_S\) 亦为自由阿贝尔群，且
\(
\rank_{\ZZ}(\ker f_S)=\mathcal L(f_S)
\)
为圆维损失（定义 \ref{def:cdim-audit-completion-loss}）。
结论由定理 \ref{thm:cdim-minimal-audit-completion} 直接给出。证毕。
\end{proof}

\begin{theorem}[部分边界 Gödelization 的最小审计补全与素数审计码]\label{thm:spg-partial-boundary-godel-audit-optimal}
沿用定理 \ref{thm:spg-partial-boundary-screen-kernel-relative-homology} 的记号，并令 \(r:=\rank_{\ZZ}(\ker f_S)\)。
取 \(\ker f_S\) 的一组 \(\ZZ\)-基 \(\{k_1,\dots,k_r\}\)。
再取互异素数 \(\{\pi_1,\dots,\pi_{2r}\}\)，并要求其与用于面标签的素数族 \(\{q_{\vec f}\}\)（定理 \ref{thm:spg-boundary-godelization-holographic-dictionary}）不交。
则存在同态
\[
\mathrm{Enc}_S:\ C_n(\mathcal Q_m;\ZZ)\longrightarrow \ZZ^{S}\times \NN,
\qquad
\mathrm{Enc}_S(x):=\bigl(f_S(x),\ \mathrm{Audit}_S(x)\bigr),
\]
满足：
\begin{enumerate}
  \item \(\mathrm{Enc}_S\) 为单射；
  \item 存在整数向量 \(\alpha(x)=(\alpha_1(x),\dots,\alpha_r(x))\in\ZZ^r\)（关于 \(x\) 的群同态），使得
        \[
        \mathrm{Audit}_S(x)
        :=
        \prod_{j=1}^{r}\pi_j^{\max(0,\alpha_j(x))}\cdot \pi_{j+r}^{\max(0,-\alpha_j(x))},
        \]
        并且由 \(\mathrm{Audit}_S(x)\) 的素因子分解可唯一恢复 \(\alpha(x)\)；
  \item 在所有使 \(f_S\) 可单射提升的审计补全中，所需的独立补账秩（圆维）最小且等于 \(r\)。
\end{enumerate}
\leanverified{paper\_spg\_partial\_boundary\_godel\_audit\_optimal}
\end{theorem}
\begin{proof}
由于 \(C_n(\mathcal Q_m;\ZZ)\) 与 \(\ZZ^{S}\) 为自由阿贝尔群，短正合列
\[
0\longrightarrow \ker f_S\longrightarrow C_n(\mathcal Q_m;\ZZ)\longrightarrow \mathrm{im}(f_S)\longrightarrow 0
\]
以 \(\mathrm{im}(f_S)\) 自由为前提可分裂，从而存在同态重traction \(p:\,C_n(\mathcal Q_m;\ZZ)\to \ker f_S\) 使 \(p|_{\ker f_S}=\mathrm{id}\)。
令 \(\alpha(x)\in\ZZ^r\) 为 \(p(x)\) 在基 \(\{k_j\}\) 下的坐标，即 \(p(x)=\sum_{j=1}^r \alpha_j(x)k_j\)。
则 \(\alpha\) 为群同态，且由算术基本定理 \(\mathrm{Audit}_S(x)\) 的素因子分解唯一确定 \(\alpha(x)\)。
\par
若 \(\mathrm{Enc}_S(x)=\mathrm{Enc}_S(y)\)，则 \(f_S(x-y)=0\) 从而 \(x-y\in\ker f_S\)，并且 \(\alpha(x-y)=\alpha(x)-\alpha(y)=0\)。
由 \(p|_{\ker f_S}=\mathrm{id}\) 得 \(p(x-y)=x-y\)，从而 \(x-y=0\)，故 \(\mathrm{Enc}_S\) 单射。
\par
最优性由推论 \ref{cor:spg-partial-boundary-min-audit-cost-kernel-rank} 给出：任意审计补全的最小成本等于 \(\rank_{\ZZ}(\ker f_S)=r\)，因此不存在秩更小的补全仍能实现单射提升。证毕。
\end{proof}

\begin{theorem}[可见秩函数的可表示拟阵性与审计缺口的超模性]\label{thm:spg-screen-rank-matroid-supermodularity}
沿用上述记号。
对任意 \(S\subset \overrightarrow{\mathcal F_{m}^{n-1}}\)，令
\[
r(S):=\rank_{\QQ}\bigl(\mathrm{im}(f_S\otimes\QQ)\bigr),
\qquad
\delta(S):=\rank_{\QQ}\bigl(\ker(f_S\otimes\QQ)\bigr).
\]
则：
\begin{enumerate}
  \item \(r(\cdot)\) 是 \(\overrightarrow{\mathcal F_{m}^{n-1}}\) 上一个可表示拟阵的秩函数；特别地，对任意 \(S\subset T\) 有 \(r(S)\le r(T)\)，并且对任意 \(S,T\) 有
        \[
        r(S)+r(T)\ \ge\ r(S\cup T)+r(S\cap T).
        \]
  \item \(\delta(S)=\rank_{\QQ}(C_n(\mathcal Q_m;\QQ)) - r(S)\) 从而满足超模性
        \[
        \delta(S)+\delta(T)\ \le\ \delta(S\cup T)+\delta(S\cap T),
        \]
        并且 \(\delta\) 随 \(S\) 单调不增。
  \item 由于 \(C_n(\mathcal Q_m;\ZZ)\) 自由，\(\rank_{\QQ}\ker(f_S\otimes\QQ)=\rank_{\ZZ}\ker f_S\)。
        结合推论 \ref{cor:spg-partial-boundary-min-audit-cost-kernel-rank}，最小审计补全成本 \(S\mapsto \min_{(R,r_S)}\cdim(R)\) 在该屏幕族上同样满足上述超模性。
\end{enumerate}
\end{theorem}
\begin{proof}
取 \(\partial_n\) 在某组 \(\ZZ\)-基下的矩阵表示 \(A\)。
则 \(f_S\otimes\QQ\) 等价于将 \(A\) 的行按 \(S\) 取子集得到的行限制矩阵 \(A_S\)，从而
\[
r(S)=\dim_{\QQ}\Span\{\text{rows of }A\text{ indexed by }S\}.
\]
由线性拟阵的一般理论，\(r(\cdot)\) 满足单调性与次模性，从而为可表示拟阵秩函数。
由秩亏恒等式得 \(\delta(S)=\rank_{\QQ}(C_n(\mathcal Q_m;\QQ))-r(S)\)，因此 \(\delta\) 为超模函数且单调不增。
最后一条由 \(\ker f_S\) 的秩在张量到 \(\QQ\) 后保持不变，以及推论 \ref{cor:spg-partial-boundary-min-audit-cost-kernel-rank} 的同一刻画直接推出。证毕。
\end{proof}

\leanverified{paper\_spg\_screen\_rank\_matroid\_supermodularity}
\leanverified{paper\_spg\_screen\_audit\_cost\_monotone\_unit\_drop\_seeds}
\leanverified{paper\_spg\_screen\_audit\_cost\_monotone\_unit\_drop\_package}
\leanverified{paper\_spg\_screen\_audit\_cost\_monotone\_unit\_drop}
\begin{corollary}[审计补全成本的单调律与单位下降界]\label{cor:spg-screen-audit-cost-monotone-unit-drop}
沿用上述记号。对任意屏幕 \(S\subseteq T\subseteq \overrightarrow{\mathcal F_m^{n-1}}\)，定义
\[
\mathfrak c(S):=\min_{(R,r_S)}\cdim(R)=\rank_{\ZZ}\ker f_S.
\]
则有单调律
\[
\mathfrak c(T)\le \mathfrak c(S).
\]
更一般地，若 \(\card(T\setminus S)=t\)，则
\[
\mathfrak c(S)-t\le \mathfrak c(T)\le \mathfrak c(S).
\]
特别地，当 \(T=S\cup\{\vec f\}\) 仅比 \(S\) 多一张取向可见面时，
\[
\mathfrak c(S)-1\le \mathfrak c(T)\le \mathfrak c(S).
\]
\end{corollary}
\begin{proof}
记
\[
d:=\rank_{\QQ}C_n(\mathcal Q_m;\QQ),
\qquad
r(S):=\rank_{\QQ}\bigl(\operatorname{im}(f_S\otimes\QQ)\bigr).
\]
由定理 \ref{thm:spg-screen-rank-matroid-supermodularity} 与推论 \ref{cor:spg-partial-boundary-min-audit-cost-kernel-rank}，
\[
\mathfrak c(S)=\rank_{\QQ}\ker(f_S\otimes\QQ)=d-r(S).
\]
当 \(S\subseteq T\) 时，\(f_T\otimes\QQ\) 由 \(f_S\otimes\QQ\) 加入若干额外行得到，故 \(r(T)\ge r(S)\)，从而
\[
\mathfrak c(T)=d-r(T)\le d-r(S)=\mathfrak c(S).
\]
若再设 \(\card(T\setminus S)=t\)，则添入 \(t\) 行至多使秩上升 \(t\)，即
\[
r(T)\le r(S)+t.
\]
于是
\[
\mathfrak c(T)=d-r(T)\ge d-r(S)-t=\mathfrak c(S)-t.
\]
当 \(t=1\) 时即得单位下降界。证毕。
\end{proof}

\paragraph{坐标束屏幕的审计缺口与外边界闭合}
在部分边界观测的屏幕核、最小审计补全成本与可表示拟阵结构已经确立之后，坐标束型内部屏幕还可给出一组完全显式的闭式计数。

\begin{theorem}[内部坐标束屏幕的连通分量与审计补全成本闭式]\label{thm:spg-internal-coordinate-bundle-screen-cost-closed-form}
固定 \(m,n\ge 1\)，记 \(\mathcal Q_m\) 为 \(I^n=[0,1]^n\) 的 \(m\)-级 dyadic \(n\)-立方复形，并以
\[
a=(a_1,\dots,a_n)\in\{0,\dots,2^m-1\}^n
\]
标记其全部顶维 \(n\)-胞腔。对任意坐标子集
\[
J\subseteq \{1,\dots,n\},
\qquad
s:=|J|,
\]
定义内部坐标束屏幕 \(S_J^{\mathrm{int}}\subset \overrightarrow{\mathcal F_m^{n-1}}\) 为全部满足下述条件的基向量所对应之取向 \((n-1)\)-面：
\begin{enumerate}
  \item 其底层几何面为内部共享面；
  \item 其法向平行于某个 \(e_j\)（\(j\in J\)）。
\end{enumerate}
则其屏幕图 \(\Gamma_{S_J^{\mathrm{int}}}\) 的连通分量数满足
\[
c(\Gamma_{S_J^{\mathrm{int}}})=2^{m(n-s)}+1,
\]
并且最小审计补全成本恰为
\[
\min_{(R,r_{S_J^{\mathrm{int}}})}\cdim(R)=2^{m(n-s)}.
\]
\end{theorem}
\begin{proof}
由定理 \ref{thm:spg-screen-kernel-connected-components} 的构造，图 \(\Gamma_{S_J^{\mathrm{int}}}\) 的顶点为全部顶维胞腔与外部顶点 \(\infty\)。其内部边恰连接那些在某个 \(j\in J\) 方向上相邻、并且在其余方向坐标完全相同的 dyadic 立方体。

因此，对任意顶维胞腔地址
\[
a=(a_1,\dots,a_n),
\]
补坐标
\[
(a_i)_{i\in J^c}
\]
在 \(\Gamma_{S_J^{\mathrm{int}}}\) 的内部子图上是不变量。固定这组补坐标以后，允许变化的仅为 \(J\) 中的 \(s\) 个坐标，故对应顶点集正好同构于
\[
\{0,\dots,2^m-1\}^s.
\]
其邻接关系是沿单个 \(J\)-坐标作一步变化所得到的网格图，即
\[
\square_{j\in J} P_{2^m},
\]
其中当 \(J=\varnothing\) 时按约定将上式解释为单点图。无论哪种情形，该图都连通。故内部顶维胞腔部分的连通分量数恰为补坐标的可能取值数
\[
(2^m)^{n-s}=2^{m(n-s)}.
\]

另一方面，由于 \(S_J^{\mathrm{int}}\) 不含任何外边界面，外部顶点 \(\infty\) 在 \(\Gamma_{S_J^{\mathrm{int}}}\) 中保持孤立。因此总连通分量数为
\[
c(\Gamma_{S_J^{\mathrm{int}}})=2^{m(n-s)}+1.
\]
最后由推论 \ref{cor:spg-screen-injectivity-and-audit-cost-components}，
\[
\min_{(R,r_{S_J^{\mathrm{int}}})}\cdim(R)
=
c(\Gamma_{S_J^{\mathrm{int}}})-1
=
2^{m(n-s)}.
\]
证毕。
\leanpartial{paper\_spg\_internal\_coordinate\_bundle\_screen\_cost\_closed\_form}{仅形式化连通分量与审计成本闭式 $2^{m(n-s)} + 1$ 与 $2^{m(n-s)}$；屏幕图 $\Gamma_{S_J^{\mathrm{int}}}$ 的连通分量论证留后轮}
\leanverified{screenComponentCount}
\leanverified{auditCost}
\leanverified{two\_pow\_m\_pow\_eq}
\leanverified{complement\_coord\_count}
\leanverified{screenComponentCount\_eq}
\leanverified{screenComponentCount\_eq\_complement\_plus\_one}
\leanverified{auditCost\_eq\_count\_sub\_one}
\leanverified{paper\_spg\_internal\_coordinate\_bundle\_screen\_cost\_closed\_form\_seeds}
\leanverified{paper\_spg\_internal\_coordinate\_bundle\_screen\_cost\_closed\_form\_package}
\end{proof}

\begin{corollary}[非空坐标束屏幕的最小外边界闭合条数]\label{cor:spg-coordinate-bundle-minimal-boundary-closure}
沿用定理 \ref{thm:spg-internal-coordinate-bundle-screen-cost-closed-form} 的记号，并额外假设 \(J\neq\varnothing\)。
设 \(B\subset \overrightarrow{\mathcal F_m^{n-1}}\) 为外边界面族。
则扩张屏幕
\[
S_J^{\mathrm{int}}\cup B
\]
为精确屏幕当且仅当 \(B\) 与 \(\Gamma_{S_J^{\mathrm{int}}}\) 的每一个内部连通分量都相接。
特别地，使其达到精确屏幕所需的外边界面最小条数恰为
\[
\min\bigl\{\card B:\ \min_{(R,r_{S_J^{\mathrm{int}}\cup B})}\cdim(R)=0\bigr\}
=
2^{m(n-s)}.
\]
\leanverified{paper\_spg\_coordinate\_bundle\_minimal\_boundary\_closure\_seeds}
\leanverified{paper\_spg\_coordinate\_bundle\_minimal\_boundary\_closure}
\end{corollary}
\begin{proof}
由定理 \ref{thm:spg-internal-coordinate-bundle-screen-cost-closed-form}，\(\Gamma_{S_J^{\mathrm{int}}}\) 的内部连通分量数恰为 \(2^{m(n-s)}\)，且外部顶点 \(\infty\) 为孤立点。
由于 \(J\neq\varnothing\)，每个内部连通分量都由一个 \(s\)-维网格块
\(
\{0,\dots,2^m-1\}^s
\)
给出；沿任意 \(j\in J\) 方向取端点坐标 \(0\) 或 \(2^m-1\)，即可得到与外边界相接的胞腔。
故每个内部连通分量都至少含有一张外边界面。

若 \(S_J^{\mathrm{int}}\cup B\) 为精确屏幕，则由推论 \ref{cor:spg-screen-injectivity-and-audit-cost-components}，其对应图必须连通。
因此每个内部连通分量都必须经由 \(B\) 中某张外边界面与 \(\infty\) 相连；否则该分量仍与 \(\infty\) 断开。
这证明了必要性。

反过来，若 \(B\) 与每个内部连通分量都相接，则每个分量都通过某条外边界边连到 \(\infty\)，从而整张图连通，故扩张屏幕精确。
又由于一张外边界面只隶属于唯一一个顶维胞腔，因而也只可能接入唯一一个内部连通分量，所以要使全部 \(2^{m(n-s)}\) 个内部连通分量都与 \(\infty\) 相连，至少需要 \(2^{m(n-s)}\) 张外边界面。
逐个连通分量各取一张即可达到该下界。证毕。
\end{proof}

\begin{corollary}[内部全坐标屏幕的唯一缺口与任意外边界闭合]\label{cor:spg-full-internal-screen-one-defect-boundary-closure}
沿用定理 \ref{thm:spg-internal-coordinate-bundle-screen-cost-closed-form} 的记号，并取
\[
[n]:=\{1,\dots,n\},
\qquad
J=[n].
\]
则内部全坐标屏幕 \(S_{[n]}^{\mathrm{int}}\) 的最小审计补全成本恰为
\[
\min_{(R,r_{S_{[n]}^{\mathrm{int}}})}\cdim(R)=1.
\]
更强地，若 \(B\subset \overrightarrow{\mathcal F_m^{n-1}}\) 是任意非空外边界面族，则扩张屏幕
\[
S_{[n]}^{\mathrm{int}}\cup B
\]
已经是精确屏幕，即
\[
\min_{(R,r_{S_{[n]}^{\mathrm{int}}\cup B})}\cdim(R)=0.
\]
\leanverified{paper\_spg\_full\_internal\_screen\_one\_defect\_boundary\_closure}
\leanverified{paper\_spg\_full\_internal\_screen\_one\_defect\_seeds}
\end{corollary}
\begin{proof}
将定理 \ref{thm:spg-internal-coordinate-bundle-screen-cost-closed-form} 应用于 \(J=[n]\) 即得
\[
c(\Gamma_{S_{[n]}^{\mathrm{int}}})=2^{m(n-n)}+1=2,
\]
从而
\[
\min_{(R,r_{S_{[n]}^{\mathrm{int}}})}\cdim(R)=c(\Gamma_{S_{[n]}^{\mathrm{int}}})-1=1.
\]

现设 \(B\) 为任意非空外边界面族。由于 \(J=[n]\)，定理
\ref{thm:spg-internal-coordinate-bundle-screen-cost-closed-form}
的证明已经表明：在内部顶维胞腔诱导的子图上，仅存在一个连通分量。另一方面，\(B\neq\varnothing\) 意味着至少存在一条外边界边把某个顶维胞腔与 \(\infty\) 相连。故在扩张图
\[
\Gamma_{S_{[n]}^{\mathrm{int}}\cup B}
\]
中，\(\infty\) 与唯一的内部连通分量相接，从而整张图连通，即
\[
c(\Gamma_{S_{[n]}^{\mathrm{int}}\cup B})=1.
\]
再由推论 \ref{cor:spg-screen-injectivity-and-audit-cost-components}，
\[
\min_{(R,r_{S_{[n]}^{\mathrm{int}}\cup B})}\cdim(R)
=
c(\Gamma_{S_{[n]}^{\mathrm{int}}\cup B})-1
=0.
\]
证毕。
\end{proof}

\begin{theorem}[坐标束屏幕的 prime-register 残差盒预算下界]\label{thm:spg-coordinate-bundle-prime-register-residue-box-budget}
沿用定理 \ref{thm:spg-internal-coordinate-bundle-screen-cost-closed-form} 的记号，并令
\[
r_J:=\min_{(R,r_{S_J^{\mathrm{int}}})}\cdim(R)=2^{m(n-s)}.
\]
取 \(\ker f_{S_J^{\mathrm{int}}}\) 的一组 \(\ZZ\)-基 \(\{k_1,\dots,k_{r_J}\}\)，并固定素数 \(p\)。
定义对应的 \(p\)-进残差盒
\[
K_{J,p}:=\left\{\sum_{i=1}^{r_J} a_i k_i:\ a_i\in\{0,\dots,p-1\}\right\}\subset \ker f_{S_J^{\mathrm{int}}}.
\]
再对整数 \(k\ge 1,\ E\ge 0\) 定义截断 prime-register 状态空间
\[
\mathcal P_{k,E}
:=
\left\{r:\mathbb P\to\NN\ \text{有限支撑}:\ r(p_i)\in\{0,\dots,E\}\ (1\le i\le k),\ r(q)=0\ (q>p_k)\right\},
\]
其中 \((p_i)_{i\ge 1}\) 为递增素数序列，从而 \(\card{\mathcal P_{k,E}}=(E+1)^k\)。
若存在映射
\[
\rho:K_{J,p}\to \mathcal P_{k,E}
\]
为单射，则必有
\[
(E+1)^k\ge p^{r_J},
\qquad\text{等价地}\qquad
k\log(E+1)\ge 2^{m(n-s)}\log p.
\]
从而：
\[
E\ge p^{\,2^{m(n-s)}/k}-1,
\qquad
k\ge \frac{2^{m(n-s)}\log p}{\log(E+1)}.
\]
此外，任何把 \(K_{J,p}\) 单射外置为单个整数的 Gödel 编码，其位长都满足
\[
\log_2 N_{\max}\ge 2^{m(n-s)}\log_2 p.
\]
\leanverified{paper\_spg\_coordinate\_bundle\_prime\_register\_residue\_box\_budget}
\end{theorem}
\begin{proof}
由 \(\{k_i\}_{i=1}^{r_J}\) 为自由基，集合 \(K_{J,p}\) 的每个元素都由 \((a_1,\dots,a_{r_J})\in\{0,\dots,p-1\}^{r_J}\) 唯一决定，故
\[
\card{K_{J,p}}=p^{r_J}.
\]
若 \(\rho:K_{J,p}\to \mathcal P_{k,E}\) 为单射，则必有
\[
\card{\mathcal P_{k,E}}\ge \card{K_{J,p}}=p^{r_J}.
\]
代入 \(\card{\mathcal P_{k,E}}=(E+1)^k\) 即得
\[
(E+1)^k\ge p^{r_J}.
\]
再取对数并用 \(r_J=2^{m(n-s)}\) 便得到
\[
k\log(E+1)\ge 2^{m(n-s)}\log p.
\]
分别解出 \(E\) 与 \(k\) 即得前两式。
若再用单个整数去单射编码 \(\card{K_{J,p}}=p^{r_J}\) 个状态，则其取值上界 \(N_{\max}\) 至少满足 \(N_{\max}\ge p^{r_J}\)，从而
\[
\log_2 N_{\max}\ge r_J\log_2 p=2^{m(n-s)}\log_2 p.
\]
证毕。
\end{proof}

\noindent
因此，坐标遗漏的代价在 dyadic 屏幕模型中并非仅呈数量级增长，而是严格满足每遗漏一个坐标方向便乘上一个 \(2^m\) 因子的闭式律；而当全部内部坐标方向均已可见时，唯一剩余缺口只是与外部参考相连的整体常值模，并且任一外边界面都足以将其消去。换言之，在定理 \ref{thm:spg-screen-rank-matroid-supermodularity} 所给出的行拟阵语言中，\(S_{[n]}^{\mathrm{int}}\) 恰处于亏秩 \(1\) 的位置，而任一外边界面都构成其达到精确屏幕的单步补足。
\begin{proposition}[坐标束闭合缺口的区间分解]\label{prop:distill-spg-coordinate-bundle-defect-intervals}
沿用定理 \ref{thm:spg-internal-coordinate-bundle-screen-cost-closed-form} 的记号，并假设 \(J\subseteq\{1,\ldots,n\}\) 非空，\(s=|J|\)。令 \(\mathcal C(J)\) 为屏幕图 \(\Gamma_{S_J^{\mathrm{int}}}\) 的内部连通分量集合；于是
\[
\card\mathcal C(J)=2^{m(n-s)}.
\]
取一条外边界审计链
\[
B_0\subseteq B_1\subseteq\cdots\subseteq B_T
\]
其中每个 \(B_t\) 为外边界面族。对 \(C\in\mathcal C(J)\)，定义
\[
d(C):=\min\{t: B_t\text{ 中某张外边界面与 }C\text{ 相接}\},
\]
若集合为空则令 \(d(C):=T+1\)。令
\[
D_t:=\bigoplus_{C\in\mathcal C(J),\ d(C)>t}\mathbb K e_C
\]
并定义转移 \(\tau_{t+1,t}:D_t\to D_{t+1}\) 为
\[
\tau_{t+1,t}(e_C)=
\begin{cases}
e_C,&d(C)>t+1,\\
0,&d(C)=t+1.
\end{cases}
\]
则 \((D_t,\tau_{t+1,t})_{0\le t\le T}\) 是一参数有限维持久模，且分解为
\[
D_\bullet\cong\bigoplus_{C\in\mathcal C(J)}\mathbb K_{[0,d(C))}.
\]
并且对每个 \(t\)，扩张屏幕 \(S_J^{\mathrm{int}}\cup B_t\) 的剩余最小外边界闭合条数恰为
\[
\dim_{\mathbb K}D_t=\#\{C\in\mathcal C(J):d(C)>t\}.
\]
\end{proposition}
\begin{proof}
由定理 \ref{thm:spg-internal-coordinate-bundle-screen-cost-closed-form}，内部坐标束屏幕的内部连通分量正由补坐标标记，数目为 \(2^{m(n-s)}\)。因为 \(J\ne\varnothing\)，推论 \ref{cor:spg-coordinate-bundle-minimal-boundary-closure} 表明每个内部连通分量均可由某张外边界面接到外部顶点 \(\infty\)，且一张外边界面只接入一个内部连通分量。

映射 \(\tau_{t+1,t}\) 显然满足 \(\tau_{u,t}=\tau_{u,u-1}\cdots\tau_{t+1,t}\)，因为每个基向量 \(e_C\) 只在首次接触尺度 \(d(C)\) 被杀死一次，之后不再出现。因此 \((D_t,\tau_{t+1,t})\) 是合法的一参数持久模。对固定 \(C\)，向量 \(e_C\) 在 \(t<d(C)\) 时存在，在 \(t\ge d(C)\) 时为零，故它贡献区间模 \(\mathbb K_{[0,d(C))}\)。对所有 \(C\) 取直和得到所示条码分解。

在尺度 \(t\)，已经接触 \(B_t\) 的内部连通分量与 \(\infty\) 同属一个连通块；未接触者仍各自与 \(\infty\) 分离。由推论 \ref{cor:spg-coordinate-bundle-minimal-boundary-closure}，要使屏幕精确，必须且只须给每个未接触分量再选一张外边界面。因此剩余闭合条数等于未死亡区间的数目，即 \(\dim_{\mathbb K}D_t\)。
\end{proof}
下述有限化使用 Interleaving stability 算子：扰动预算不比较外边界面本身，而比较其在坐标束屏幕图连通分量载体上的首次闭合时间。
\begin{theorem}[坐标束闭合审计的载体死亡时间稳定定理]\label{distill:persistent_homology_persistence_modules_and_constructible_sheaf_persistence-stability_and_interleaving-spg-coordinate-bundle-carrier-death-stability}
沿用定理 \ref{thm:spg-internal-coordinate-bundle-screen-cost-closed-form} 与推论 \ref{cor:spg-coordinate-bundle-minimal-boundary-closure} 的记号，并假设 \(J\neq\varnothing\)。置
\[
\Lambda_J:=\{0,\dots,2^m-1\}^{J^c}.
\]
若 \(b\) 是一个外边界取向面，记 \(Q_a\) 为与 \(b\) 相接的唯一顶维 dyadic 立方体，并定义其坐标束载体为
\[
\kappa(b):=(a_i)_{i\in J^c}\in \Lambda_J.
\]
设 \(L\ge 1\)，称递增外边界面族
\[
B_0\subseteq B_1\subseteq\cdots\subseteq B_L
\]
为闭合有限外边界审计协议，如果对每个 \(u\in\Lambda_J\) 都存在 \(b\in B_L\) 满足 \(\kappa(b)=u\)。定义载体 \(u\) 的首次闭合时间
\[
\tau_B(u):=\min\{t\in\{0,\dots,L\}:\exists b\in B_t\text{ with }\kappa(b)=u\}.
\]
令 \(\mathbb F\) 为任意域，并定义有限审计 ledger
\[
M_B(t):=\operatorname{span}_{\mathbb F}\{e_u:u\in\Lambda_J,\ t<\tau_B(u)\},\qquad 0\le t\le L,
\]
其中 \(M_B(t)\to M_B(t')\) 对 \(t\le t'\) 为删去已死亡基向量的坐标投影。

若两个闭合审计协议 \(B\) 与 \(C\) 满足如下 \(\rho\)-载体交错条件：对所有 \(0\le t\le L\)，
\[
\{u:\tau_B(u)\le t\}\subseteq \{u:\tau_C(u)\le \min(L,t+\rho)\},
\]
且同一包含式在交换 \(B,C\) 后仍成立，则
\[
|\tau_B(u)-\tau_C(u)|\le \rho
\qquad\text{for every }u\in\Lambda_J.
\]
特别地，若
\[
d_{\mathrm{death}}(B,C):=\min_{\sigma\in\operatorname{Sym}(\Lambda_J)}\max_{u\in\Lambda_J}|\tau_B(u)-\tau_C(\sigma(u))|,
\]
则
\[
d_{\mathrm{death}}(B,C)\le \rho.
\]
此外，对任意截止深度 \(r\in\{0,\dots,L\}\)，两个审计 ledger 在 \(r\) 处的生存判决只能在 \(B\)-死亡时间的 \(\rho\)-邻域内不同：
\[
\{u:(e_u\in M_B(r))\not\Longleftrightarrow(e_u\in M_C(r))\}
\subseteq
\{u:r-\rho<\tau_B(u)\le r+\rho\}.
\]
因此，若某载体 \(u\) 在 \(B\) 中满足 \(\tau_B(u)>r+\rho\)，则它在所有与 \(B\) \(\rho\)-载体交错的闭合协议 \(C\) 中仍满足 \(e_u\in M_C(r)\)；若 \(\tau_B(u)\le r-\rho\)，则在所有此类 \(C\) 中均有 \(e_u\notin M_C(r)\)。
\end{theorem}
\begin{proof}
由定理 \ref{thm:spg-internal-coordinate-bundle-screen-cost-closed-form} 的证明，\(S_J^{\mathrm{int}}\) 的内部连通分量正由补坐标 \(u\in\Lambda_J\) 标记。由于 \(J\neq\varnothing\)，固定补坐标后，允许变化的 \(J\)-坐标形成一个非空坐标方向上的有限网格；该网格与外边界相接。因此一个外边界面 \(b\) 的载体 \(\kappa(b)\) 正好指定它闭合的内部连通分量。

闭合协议的假设给出 \(\tau_B(u),\tau_C(u)\in\{0,\dots,L\}\)。取任意 \(u\in\Lambda_J\)，并在 \(\rho\)-载体交错条件中令 \(t=\tau_B(u)\)。因 \(u\in\{v:\tau_B(v)\le t\}\)，得到
\[
\tau_C(u)\le \min(L,\tau_B(u)+\rho)\le \tau_B(u)+\rho.
\]
交换 \(B,C\) 后同理得到
\[
\tau_B(u)\le \tau_C(u)+\rho.
\]
两式合并即为 \(|\tau_B(u)-\tau_C(u)|\le\rho\)。取置换 \(\sigma=\operatorname{id}_{\Lambda_J}\) 便得 \(d_{\mathrm{death}}(B,C)\le\rho\)。

再固定截止深度 \(r\)。按定义，\(e_u\in M_B(r)\) 当且仅当 \(r<\tau_B(u)\)，而 \(e_u\notin M_B(r)\) 当且仅当 \(\tau_B(u)\le r\)。若 \(\tau_B(u)\le r-\rho\)，则由上段不等式
\[
\tau_C(u)\le \tau_B(u)+\rho\le r,
\]
故 \(e_u\notin M_C(r)\)，不会发生生存判决差异。若 \(\tau_B(u)>r+\rho\)，则
\[
\tau_C(u)\ge \tau_B(u)-\rho>r,
\]
故 \(e_u\in M_C(r)\)，同样不会发生差异。于是任何差异载体必须满足
\[
r-\rho<\tau_B(u)\le r+\rho.
\]
最后两项断言正是这两个排除情形的重述。推论 \ref{cor:spg-coordinate-bundle-minimal-boundary-closure} 保证这里的每个载体死亡都对应一个内部连通分量被外边界闭合，故该死亡时间判据正是坐标束屏幕审计闭合的有限稳定判据。证毕。
\end{proof}
\begin{definition}[屏幕图的有限 Dirac 距离预算]\label{def:distill-connes-screen-dirac-distance-budget}
设 \(\Gamma=(V,E)\) 为有限无向图，且每条边 \(e=\{v,w\}\) 配有权重 \(\kappa_e>0\)。令
\[
H_{\Gamma}:=\ell^2(V),\qquad \mathcal A_{\Gamma}:=\RR^V,
\]
并令 \(M_f\) 表示 \(f\in\mathcal A_{\Gamma}\) 的乘法算子。定义自伴更新算子
\[
D_{\kappa}\delta_v:=\sum_{w:\{v,w\}\in E}\kappa_{\{v,w\}}\delta_w.
\]
置
\[
L_{\kappa}(f):=\|[D_{\kappa},M_f]\|,\qquad
\mathrm d_{\kappa}(u,v):=\sup\{|f(u)-f(v)|:\ f\in\mathcal A_{\Gamma},\ L_{\kappa}(f)\le1\}.
\]
若上确界无界，则约定 \(\mathrm d_{\kappa}(u,v)=\infty\)。
\end{definition}

\begin{theorem}[屏幕闭合的交换子距离判据]\label{thm:distill-connes-screen-commutator-distance-closure}
令 \(S\) 为任一屏幕面族，\(\Gamma_S\) 为前文定义的屏幕图，并在 \(\Gamma_S\) 的每条边上取正权重 \(\kappa\)。则以下命题成立：
\begin{enumerate}
  \item 若顶点 \(u,v\) 属于 \(\Gamma_S\) 的不同连通分量，则 \(\mathrm d_{\kappa}(u,v)=\infty\)。
  \item 若 \(u=v_0,v_1,\ldots,v_r=v\) 是 \(\Gamma_S\) 中的一条路径，则
  \[
  \mathrm d_{\kappa}(u,v)\le \sum_{i=0}^{r-1}\kappa_{\{v_i,v_{i+1}\}}^{-1}.
  \]
  特别地，同一连通分量内的任意两点具有有限 \(\mathrm d_{\kappa}\)-距离。
\end{enumerate}
沿用定理 \ref{thm:spg-internal-coordinate-bundle-screen-cost-closed-form} 的记号。若 \(J\ne\varnothing\)，则对扩张屏幕 \(S_J^{\mathrm{int}}\cup B\)，所有内部顶维胞腔到外部顶点 \(\infty\) 的 \(\mathrm d_{\kappa}\)-距离均有限，当且仅当 \(B\) 与 \(\Gamma_{S_J^{\mathrm{int}}}\) 的每个内部连通分量都相接。因此达到该性质所需的外边界面最小条数为 \(2^{m(n-s)}\)。
\end{theorem}
\begin{proof}
若 \(u,v\) 分属不同连通分量，令 \(C\) 为含 \(u\) 的连通分量，并令 \(f=t\ind_C\)，其中 \(t>0\)。由于 \(\Gamma_S\) 没有连接 \(C\) 与其补集的边，对每条边 \(\{x,y\}\) 都有 \(f(x)=f(y)\)。故 \([D_{\kappa},M_f]=0\)，即 \(L_{\kappa}(f)=0\le1\)，而 \(|f(u)-f(v)|=t\)。令 \(t\to\infty\) 得 \(\mathrm d_{\kappa}(u,v)=\infty\)。

若 \(\{x,y\}\) 为一条边且 \(L_{\kappa}(f)\le1\)，则
\[
[D_{\kappa},M_f]\delta_y
=\sum_{z:\{y,z\}\in E}\kappa_{\{y,z\}}(f(y)-f(z))\delta_z,
\]
因此
\[
1\ge \|[D_{\kappa},M_f]\delta_y\|\ge
\kappa_{\{x,y\}}|f(y)-f(x)|.
\]
于是 \(|f(x)-f(y)|\le\kappa_{\{x,y\}}^{-1}\)。沿路径求和得到
\[
|f(u)-f(v)|\le \sum_{i=0}^{r-1}\kappa_{\{v_i,v_{i+1}\}}^{-1}.
\]
对所有满足 \(L_{\kappa}(f)\le1\) 的 \(f\) 取上确界，即得第二项。

对 \(S_J^{\mathrm{int}}\cup B\)，顶维胞腔到 \(\infty\) 的谱距离有限正等价于它们在扩张屏幕图中与 \(\infty\) 同属一个连通分量。由推论 \ref{cor:spg-coordinate-bundle-minimal-boundary-closure}，这又等价于 \(B\) 与 \(\Gamma_{S_J^{\mathrm{int}}}\) 的每个内部连通分量都相接，并且最小条数为 \(2^{m(n-s)}\)。证毕。
\end{proof}

\endinput


% --- Distillation writeback ---
\begin{proposition}[部分屏幕核的相位 holonomy 寄存器下界]\label{distill:bourgain-multiscale_decoupling_gates-spg-screen-phase-holonomy-register-lower-bound}
沿用定理 \ref{thm:spg-partial-boundary-screen-kernel-relative-homology}、定理 \ref{thm:spg-screen-kernel-connected-components} 与推论 \ref{cor:spg-screen-injectivity-and-audit-cost-components} 的记号。固定屏幕 \(S\)，并记
\[
K_S:=\ker f_S\cong H_n(\mathcal Q_m,A_{\neg S};\ZZ).
\]
令 \(D\ge2\)，并设
\[
\mu_D:=\{z\in\mathbb C:z^D=1\}.
\]
给定一个 screen-invisible 相位 holonomy 同态
\[
h:K_S/DK_S\longrightarrow H_D,
\]
其中 \(H_D\) 为有限阿贝尔群；典型情形为 \(H_D\le(\mu_D)^d\)。称有限寄存器 \(R\) 与映射
\[
r:C_n(\mathcal Q_m;\ZZ)\longrightarrow R
\]
对 \((S,h)\) 相位完备，若对任意 \(u,v\in C_n(\mathcal Q_m;\ZZ)\)，条件
\[
f_S(u)=f_S(v),
\qquad
r(u)=r(v)
\]
蕴含
\[
h([u-v])=1,
\]
其中 \([u-v]\) 表示 \(u-v\in K_S\) 在 \(K_S/DK_S\) 中的类。则任意相位完备寄存器满足
\[
|R|\ge |\operatorname{im}h|.
\]
特别地，若 \(D=p\) 为素数且 \(H_D=(\mu_p)^d\)，把 \(h\) 视为 \(\mathbb F_p\)-线性映射
\[
K_S/pK_S\longrightarrow (\mathbb F_p)^d
\]
并记其秩为 \(r_h\)，则
\[
|R|\ge p^{r_h}.
\]
当 \(H_D=\mu_p\) 为单个循环相位轴时，该结论只给出非平凡情形下的下界 \(|R|\ge p\)，而不把整个 screen kernel 秩误认为单相位轴的寄存器维数。
\end{proposition}
\begin{proof}
取任意 \(u_0\in C_n(\mathcal Q_m;\ZZ)\)。对每个 \(k\in K_S\)，有
\[
f_S(u_0+k)=f_S(u_0),
\]
因为 \(K_S=\ker f_S\)。选择 \(K_S/DK_S\) 中关于 \(\ker h\) 的商
\[
(K_S/DK_S)/\ker h
\]
的一组代表元 \(\bar k_1,
\ldots,
\bar k_M\)，其中
\[
M=|(K_S/DK_S)/\ker h|=|\operatorname{im}h|,
\]
并取其在 \(K_S\) 中的任意提升 \(k_1,
\ldots,
 k_M\)。

若存在 \(a\ne b\) 使
\[
r(u_0+k_a)=r(u_0+k_b),
\]
则由于二者有相同的部分屏幕读出，相位完备性给出
\[
h([k_a-k_b])=1.
\]
这意味着 \(\bar k_a\) 与 \(\bar k_b\) 在 \((K_S/DK_S)/\ker h\) 中为同一类，和代表元选择矛盾。因此
\[
u_0+k_1,
\ldots,
u_0+k_M
\]
在寄存器 \(R\) 中有两两不同的像，故 \(|R|\ge M=|\operatorname{im}h|\)。

若 \(D=p\) 且 \(H_D=(\mu_p)^d\)，经固定 \(\mu_p\cong\mathbb F_p\) 的群同构后，\(h\) 是有限维 \(\mathbb F_p\)-线性映射，其像大小为 \(p^{r_h}\)。代入前述一般下界即得 \(|R|\ge p^{r_h}\)。证毕。
\end{proof}

% --- End distillation writeback ---

\endinput
